\section*{Chapter \ref{ch:nw:servers} --- Servers}

\begin{solution}{exo:nw:servers:1}
$2\,000 \times (0.8 \times 3/5 + 0.2) = \qty{1\,360}{\nano\second}$; the limit is the memory-bound part, $0.2 \times 2\,000 =
\qty{400}{\nano\second}$.
\end{solution}

\begin{solution}{exo:nw:servers:2}
$15\,000/570 = 26.3$: 26 servers (space allows 42). $15\,000/1\,250 = 12$ servers of 2 units (space allows 21): power binds in both.
\end{solution}

\begin{solution}{exo:nw:servers:3}
The boost, if the hot path runs on one or two cores and the rest of the chip is kept quiet (its thermal and power headroom is what lets one
core boost); the base clock if many cores are busy at once, which is when the processor cannot hold the boost on all of them.
\end{solution}

\begin{solution}{exo:nw:servers:4}
SMIs are handled by firmware outside the operating system and cannot be masked, so no operating-system setting moves them off a core. The firm
can ask the vendor which SMI sources are safe to disable, choose firmware with fewer of them, count them with the processor's counters, and
correlate unexplained pauses with them.
\end{solution}

\begin{solution}{exo:nw:servers:5}
\qty{367}{\nano\second} if all core-bound, \qty{220}{\nano\second} at 60\%, \qty{110}{\nano\second} at 30\%.
\end{solution}

\begin{solution}{exo:nw:servers:6}
A hot path that uses several cores at once (a feed handler, a strategy, a gateway on separate cores) runs at the base clock of a normal part
when they are all busy, and a load-dependent clock moves its latency with the load. An all-core overclock runs every core at one fixed
frequency, so the path's latency does not depend on what else is running.
\end{solution}

\begin{solution}{exo:nw:servers:7}
Below 25\% core-bound: $1 - (0.25 \times 4.8/5.0 + 0.75) = 1\%$. At \qty{20}{\kilo\watt}: 35 of the 9175F servers (560 cores), 38 of the 6544Y
(608), 16 of the two-socket 9755 (4\,096 cores), 10 overclocked servers (240).
\end{solution}

\begin{solution}{exo:nw:servers:8}
If the path is mostly memory- or device-bound, a faster clock changes little (\cref{prop:nw:servers:clock}); and the new server may not have
been boosting (other busy cores, a power profile, deep idle states), or the path may run on the wrong socket. Measure the core-bound share and
the clock the core actually ran at before blaming the benchmark.
\end{solution}

\begin{solution}{pb:nw:servers:1}
\begin{enumerate}
\item \qty{1\,491}{\nano\second} (9175F at 5.0~GHz), \qty{1\,685}{\nano\second} (9755 at 4.1), \qty{1\,528}{\nano\second} (overclocked, 4.8).
\item \qty{1\,659}{\nano\second}, \qty{2\,241}{\nano\second} and \qty{1\,528}{\nano\second}.
\item $2\,241 / 1\,491 = 1.50$.
\item $1/(1 - 0.6) = 2.5$: at most the memory-bound \qty{611}{\nano\second} would remain.
\item 17, 8 and 5.
\item 272, 2\,048 and 120.
\item Power for all three (space would allow 42, 21 and 21).
\item Three 9175F servers (48 hot cores) and four 9755 servers (1\,024 cores): $3 \times 570 + 4 \times 1\,250 = 6\,710$~W, one cabinet.
\item The power profile and idle states (a machine that sleeps is not the one that was tested), the turbo and frequency lock (the reason it was
  bought), memory speed and ECC reporting (overclocked memory), and the SMI sources.
\item Failures that appear only under sustained load and heat: marginal stability of the overclock, memory errors, throttling.
\item Check the fans and the liquid loop, compare with the burn-in's envelope, move its strategies to a spare before it throttles or fails,
  and open a case with the vendor.
\item A sibling thread shares the core's pipeline and caches with the hot thread, adding contention and jitter; leaving it idle keeps the core
  to the hot thread.
\item \textbf{Named result.} At 60\% core-bound the hot path takes 1\,491, 1\,685 and \qty{1\,528}{\nano\second} on the 9175F, the 9755 and the
  overclocked server; a \qty{10}{\kilo\watt} cabinet holds 17, 8 and 5 of them.
\item Below about 6\% core-bound, even the slowest clock (2.7~GHz) is within 5\% of the fastest.
\item In research, simulation and the firm's other work (One Quant Book~15), in a data centre with cheap power, not in the colocation cage.
\item The path on two clocks (or two cores of different clocks), fitted to \cref{prop:nw:servers:clock}, and the clock the core actually ran at.
\item A clock that holds still keeps the latency distribution tight; a higher but moving clock adds variance at the times the machine is
  busiest.
\item Failures are the vendor's to fix, on its terms: spares, on-site support and replacement times go into the contract (chapter~15).
\item ``The 16-core 5~GHz servers give the hot path its best clock at the lowest power; an overclocked machine is worth its premium only if our
  hot path uses many cores at once. The 128-core servers belong in research, not in the cage.''
\item The share of the path's time that is core-bound.
\end{enumerate}
\end{solution}

\begin{solution}{iq:nw:servers:1}
A hot path runs on a few cores; what matters is how fast each of them runs, and a part with fewer cores can boost higher and keeps more cache
and power per core. Many cores add throughput the hot path does not use, and lower clocks.
\iqlookfor{per-core speed against throughput.}
\end{solution}

\begin{solution}{iq:nw:servers:2}
Maximum-performance power profile, deep C-states off, a fixed turbo policy, SMT off or unused on hot cores, memory at rated speed without
cross-socket interleaving, error reporting and SMI sources minimised as the vendor allows; then audit them with the OS settings.
\iqlookfor{each setting with the latency mechanism it addresses.}
\end{solution}

\begin{solution}{iq:nw:servers:3}
An interrupt handled by firmware, invisible to and not maskable by the OS, that stops the processor for its duration (microseconds to
milliseconds); core isolation cannot remove it, so it shows up as unexplained pauses.
\iqlookfor{firmware, not maskable, not isolatable.}
\end{solution}

\begin{solution}{iq:nw:servers:4}
If the hot path is core-bound and uses several cores at once. Ask for the frequency under our sustained load, the stability testing and burn-in,
the cooling's failure modes, power draw, the warranty and spares, the firmware settings and SMI sources, and references from comparable users.
\iqlookfor{sustained behaviour and support, not the peak number.}
\end{solution}

\begin{solution}{iq:nw:servers:5}
Measure the path's core-bound share by running it at two clocks and fitting $t(f) = t_{\mathrm{ref}}(\phi f_{\mathrm{ref}}/f + 1 - \phi)$, then
predict; confirm on a loan server with the firmware configured as in production.
\iqlookfor{a model from measurement, then a test.}
\end{solution}

\begin{solution}{iq:nw:servers:6}
Power first: budget each server at its real draw and fill the cabinet with hot-path servers, switches, timing and capture; everything that does
not need to be close to the venue goes elsewhere. Keep a spare and headroom for growth.
\iqlookfor{power as the binding constraint and a clear rule for what belongs in the cage.}
\end{solution}
